\chapter{Background and related work}
\label{cha:1}
{\color{red} A chapter is a logical unit. It normally starts with an introduction, which
you are reading now. The last topic of the chapter holds the conclusion.}

\section{Current Fixture Planning Approaches}
First comes the introduction to this topic.

\subsection{Rule Based Generative fixture design}


\begin{itemize}
	\item yes no decision trees
	\item hardcoded geometric constraints
	\item deterministic logic means low flexibility
	\item brittle with more difficult geometry
	\item low adaptability
\end{itemize}

\begin{figure}
	\centering
	\includegraphics[scale=0.7]{images/Decision tree generative.png}
	\caption{Generative fixture design decision tree (source)}
	\label{fig:generative}
\end{figure}

\subsection{Case based variant fixture design}
\begin{itemize}
	\item Group technology classification (eg opitz)
	\item case library of prior part solutions
	\item similarity search
	\item adaptation step
	\item dataset quality determines performance
	\item limited innovation because biased toward historical design, unless adaptation step is very sophisticated
	\item adaptation hard for new features
	\item schematic \ref{fig:variant}
	 
	 
	 \begin{figure}
	 	\centering
	 	\includegraphics[scale=0.7]{images/variant.png}
	 	\caption{Variant fixture design scheme (source)}
	 	\label{fig:variant}
	 \end{figure}
	 
	 
	 
\end{itemize}

\section{AI in manufacturing}
\subsection{General role of AI in manufacturing}
\begin{itemize}
	\item data driven optimization
	\item reducing human dependency
\end{itemize}
\subsection{AI for CAD/CAM applications}
\begin{itemize}
	\item feature/pattern recognition
	\item automatic toolpath generation
	\item geometric classification (paper from presentation) \ref{fig:deeplearningclassification}
\end{itemize}

 \begin{figure}
	\centering
	\includegraphics[scale=0.4]{images/deeplearningclassification.png}
	\caption{Deep learning classification}
	\label{fig:deeplearningclassification}
\end{figure}


\subsection{Motivation for AI based methods}
\begin{itemize}
	\item abstraction beyond hard coded rules
	\item ability to handle many geometric inputs without breaking (robustness)
	\item probabilistic prediction
	\end{itemize}

\subsection{Challenges \& Limitations}
\begin{itemize}
	\item large labeled dataset needed
	\item computational cost
	\item explainibility issues (black box)
	\item safety concerns with decision making (liability?)
\end{itemize}
\subsection{Trends in Industry}
\begin{itemize}
	\item Move toward hybrid human-AI system, rather than fully autonomous decision making
	\item digital twins and AI monitoring
	\item predictive maintenance
\end{itemize}

\section{Conclusion}
\begin{itemize}
	\item classical methods rigid, brittle and have limited adaptability
	\item ai enables datadriven reasoning, robustness, abstraction, probablistic outputs
	\item but has its challenges such as data requirements, computational cost, explainability, safety
	\item industry trend: hybrid workflow
\end{itemize}